\section{Ejemplos}\label{sec:ejemplos}

\begin{frame}
\frametitle{Ejemplos}
Mostrar algunos ejemplos básicos, para que la gente tenga una idea del código que produce distintas cosas, uso de paquetes.

\

Indicar dónde se puede obtener más información sobre cómo escribir cosas con \LaTeX, \href{https://ctan.org/}{CTAN},  documentación de \href{https://www.overleaf.com/learn}{Overleaf}, \href{https://tex.stackexchange.com/}{\TeX{} Stack Exchange}, \href{https://www.latex-project.org/help/documentation/}{The \LaTeX{}   Project}
    
\end{frame}

\begin{frame}[fragile]
\frametitle{Ejemplos}
\framesubtitle{Hola mundo!}

Como es habitual con muchos lenguajes, veamos cómo hacer nuestro "Hola mundo!":

\begin{multicols}{2}
\begin{lstlisting}[title={hola\_mundo.tex}]
\documentclass{article}

\begin{document}

Hola mundo!

\end{document}
\end{lstlisting}

\begin{figure}
hola\_mundo.pdf
\includegraphics[width=0.35\textwidth]{../images/ejemplo_hola_mundo.png}
\end{figure}

\end{multicols}

\end{frame}